% Part: normal-modal-logic
% Chapter: frame-correspondence
% Section: equivalence-S5

\documentclass[../../../include/open-logic-section]{subfiles}

\begin{document}

\olfileid{nml}{frd}{es5}

\olsection{等价关系与 \Log{S5}}

模态逻辑 \Log{S5} 被刻画为在所有全通框架（即每个世界都能从每个世界通达，包括从自身通达自身的框架）上有效的公式的集合。在此情形下，$\Box$ 对应于必然性而 $\Diamond$ 对应于可能性：$\Box !A$ 为真当且仅当 $!A$ 在\emph{所有}世界中为真，$\Diamond !A$ 为真当且仅当 $!A$ 在\emph{某个}世界中为真。事实上，\Log{S5} 也可以被刻画为在所有自反、对称且传递的框架上有效的公式，即可及关系为\emph{等价关系}的所有框架。

\begin{defn}
  $W$ 上的二元关系 $R$ 是一个\emph{等价关系}，当且仅当它是自反、对称且传递的。$W$ 上的关系 $R$ 是\emph{全通的}，当且仅当对所有 $u,v \in W$ 都有 $Ruv$。
\end{defn}

由于 \Ax{T}、\Ax{B} 和 \Ax{4} 分别刻画了自反、对称和传递的框架，因此可及关系是等价关系的框架恰好就是这三个公式皆有效的那些框架。事实上，等价关系也可以用其他公式组合来刻画，因为我们用来定义等价关系的条件，等价于 $R$ 的其他常见条件的组合。

\begin{prop}\ollabel{prop:equivalences}
  下列条件等价：
  \begin{enumerate}
  \item $R$ 是等价关系；
  \item $R$ 是自反且欧几里得的；
  \item $R$ 是持续的、对称的且欧几里得的；
  \item $R$ 是持续的、对称的且传递的。
  \end{enumerate}
\end{prop}

\begin{proof}
  练习。
\end{proof}

\begin{prob}
  通过证明以下各条来证明 \olref[nml][frd][es5]{prop:equivalences}：
  \begin{enumerate}
  \item 若 $R$ 对称且传递，则它是欧几里得的。
  \item 若 $R$ 自反，则它是持续的。
  \item 若 $R$ 自反且欧几里得，则它是对称的。
  \item 若 $R$ 对称且欧几里得，则它是传递的。
  \item 若 $R$ 持续、对称且传递，则它是自反的。
  \end{enumerate}
  解释为什么这足以证明这些条件是等价的。
\end{prob}

\olref{prop:equivalences} 是 \olref[prf][prs]{prop:S5} 的语义对应物，因为它给出了可及关系 $R$ 为等价关系的框架的模态逻辑（传统上称为 \Log{S5}）的一个等价刻画。

全通关系与等价关系之间有何联系？尽管每个全通关系都是等价关系，但显然并非每个等价关系都是全通的。然而，在所有全通关系上有效的公式，与在所有等价关系上有效的公式完全相同。

\begin{prop}
  设 $R$ 是等价关系，并对每个 $w \in W$ 定义 $w$ 的\emph{等价类}为集合 $[w] = \{w'\in W :
  Rww'\}$。则：
  \begin{enumerate}
  \item $w \in [w]$；
  \item $R$ 在每个等价类 $[w]$ 上是全通的；
  \item 所有等价类的集合将 $W$ 划分为互不相交且穷尽 $W$ 的子集。
  \end{enumerate}
\end{prop}

\begin{prop}\ollabel{prop:S5=univ}
  公式~$!A$ 在所有 $R$ 为等价关系的框架 $\mModel{F} =\tuple{W,R}$ 上有效，当且仅当它在所有 $R$ 为全通关系的框架 $\mModel{F} =\tuple{W,R}$ 上有效。
  因此，全通框架的逻辑正是 \Log{S5}。
\end{prop}

\begin{proof}
  显然，$W$ 上的全通关系 $R$ 是等价关系。因此，若 $!A$ 在所有 $R$ 是等价关系的框架上有效，则它在所有全通框架上也有效。对于另一个方向，我们采用逆否法证明：假设公式 $!B$ 在基于框架 $\tuple{W,R}$（其中 $R$ 是 $W$ 上的等价关系）的模型 $\mModel{M} = \tuple{W, R, V}$ 中的世界 $w$ 处不成立。所以 $\mSat/{M}{!B}[w]$。定义模型 $\mModel{M}' = \tuple{W',
    R', V'}$ 如下：
  \begin{enumerate}
  \item $W' = [w]$；
  \item $R'$ 在 $W'$ 上是全通的；
  \item $V'(p) = V(p) \cap W'$。
  \end{enumerate}
  （因此，$\mModel{M}'$ 中的世界集合 $W'$ 由 \olref{fig:partition} 中的阴影区域表示。）容易看出 $R$ 和 $R'$ 在 $W'$ 上是一致的。然后可以通过对公式进行归纳证明，对所有 $w' \in W'$ 和每个公式 $!A$：$\mSat{M'}{!A}[w']$ 当且仅当 $\mSat{M}{!A}[w']$ （这是有意义的，因为 $W' \subseteq W$）。特别地，$\mSat/{M'}{!B}[w]$，因此 $!B$ 在一个基于全通框架的模型中不成立。
\end{proof}

\begin{figure}[t]
  \centering
  \begin{tikzpicture}[node distance=2cm, auto, thick]
    \clip [rounded corners] (0,0) -- (6,0) -- (6,4) -- (0,4) --  cycle;
    \begin{scope}
      \clip (0,0) -- (2,0)  .. controls (1.5,1.5) and (4.5,1.5)
      .. (4,4) -- (0,4) -- (0,0) -- cycle;
      \filldraw[fill=gray!40] (0,4) -- (0,2) .. controls (1.5,1.5)
      and (4.5,1.5) .. (4.5,0) -- (6,0) -- (6,4) -- (0,4) --  cycle;
    \end{scope}
    \draw [name path=line1] (2,0) .. controls (1.5,1.5) and (4.5,1.5) .. (4,4);
    \draw [name path=line2] (0,2)  .. controls (1.5,1.5) and (4.5,1.5) .. (4.5,0);
    \path [name intersections={of=line1 and line2}];
    \draw [rounded corners, use as bounding box, very thick] (0,0)
    -- (6,0) -- (6,4) -- (0,4) --  cycle; 
    \node at (2,3) {$[w]$} ;
    \node at (1,1) {$[u]$} ;
    \node at (3,0.75) {$[v]$} ;
    \node at (4.75,2) {$[z]$} ;
  \end{tikzpicture}
  \caption{$W$ 按等价类划分示意图。}
  \ollabel{fig:partition}
\end{figure}

\end{document}